\section{The binary model and the padded graph}\label{sec:graph}

This section separates bounded small coordinates from a complete rectangle
of large coordinates. Its repair map controls the mass of defect states;
padding then gives a uniform lower bound on the chance of recovering an
original table from a stationary state and its completion.

We count bit operations and use independent unbiased random bits.
All margins are nonnegative integers given in binary, with common total
$N$.  Accuracy is specified by an integer $k\ge1$, meaning error at most
$2^{-k}$.
An empty dimension forces $N=0$ and gives a unique empty table.
If $N=0$, or if one dimension is one, the unique table can likewise be
returned directly.  In all remaining constructions the dimensions are
positive, and we use
\begin{equation}\label{eq:parameters}
\begin{gathered}
 d=10+(m+1)(n+1),\qquad L=d^{12},\qquad U=d^{20},\\
 C=N+dL+U+2,\qquad b=d+\lceil\log_2(C+2)\rceil.
\end{gathered}
\end{equation}
Here $d$ is a dimension allowance, $L$ the padding, $U$ the small-margin
threshold, $C$ a coordinate cap, and $b$ a bit-size allowance.
Thus $d\ge14$, $mn<d$, $U<C$, and
$b=O(d+\log(N+1))$.  Polynomial bounds in $d,b,k$ have the required
binary-input dependence.  A fixed feasible table is obtained greedily:
choose a row and column of positive remaining margins, place the minimum
of those two margins in their cell, and subtract it from both margins.
Each placement exhausts a row or column.  This uses at most $m+n-1$
positive placements and polynomial bit work, including initialization of
the output.  We fix index order for this construction and all later ties.

\subsection{Small states and large completions}

Put $I=\{i:r_i\ge U\}$ and $J=\{j:c_j\ge U\}$.
The \emph{large slots} are the cells in the complete rectangle
$\mathcal B=I\times J$; the remaining cells form the set $\mathcal S$ of
\emph{small slots}.  Set $B_s=C$ on $\mathcal B$ and $B_s=U$ on
$\mathcal S$.  Associate to each slot two coordinates
$x_s,y_s\in\{0,\ldots,B_s\}$, and put $q_s=B_s-y_s$.
Here $x$ is the row viewpoint and $q$ is the column viewpoint.
A slot is \emph{balanced} when $x_s=q_s$, equivalently
$x_s+y_s=B_s$.
The balanced representation of a value $a\in\{0,\ldots,B_s\}$ is
$(x_s,y_s)=(a,B_s-a)$; both $x_s$ and $q_s$ then equal $a$.

A graph state retains only the coordinates on $\mathcal S$.
Its allowed pattern is either a \emph{transversal}, with all slots
balanced, or a \emph{defect}, with distinct slots $t,s$ satisfying
\[
 x_t=q_t+1,\qquad x_s=q_s-1,
 \qquad x_v=q_v\quad(v\notin\{s,t\}).
\]
We call $t$ positive and $s$ negative.  Both patterns have total doubled
occupancy $\sum_{s\in\mathcal S}(x_s+y_s)=U|\mathcal S|$.
Define the original residual margins by
\[
 R_i^0(X)=r_i-\sum_{j:(i,j)\in\mathcal S}x_{ij},\qquad
 P_j^0(X)=c_j-\sum_{i:(i,j)\in\mathcal S}q_{ij}.
\]
Their totals agree because $\sum_{s\in\mathcal S}(x_s-q_s)=0$.
Let $k_i$ and $\ell_j$ be the numbers of large slots in row $i$ and
column $j$.  A state is \emph{feasible} when all residuals are nonnegative
and $R_i^0=0$ if $k_i=0$, $P_j^0=0$ if $\ell_j=0$.
Its padded residuals and enlarged full margins are
\begin{align*}
 R_i(X)&=R_i^0(X)+k_iL,& P_j(X)&=P_j^0(X)+\ell_jL,\\
 \bar r_i&=r_i+k_iL,&\bar c_j&=c_j+\ell_jL.
\end{align*}
Write $\mathcal F_X$ for the set of nonnegative integer tables on
$\mathcal B$ with margins $R_i(X),P_j(X)$, and set
$f(X)=|\mathcal F_X|$.  On infeasible states set $f(X)=0$.
The complete rectangle admits a greedy completion whenever these
residual tests hold.  If $\mathcal B$ is empty the only completion is
the empty table, so feasible states have weight one.

Every completion has total at most $N+|\mathcal B|L<C$, and every
nonempty block margin is at least $L$.  In particular the formal large
capacity $C$ excludes no completion: its doubled coordinates can be
taken as $x_s=A_s,y_s=C-A_s$ for $A\in\mathcal F_X$.
The bounded sampler will draw completions of this block without
enumerating its large coordinate ranges.  One-row or one-column blocks have a
unique completion and need no random draw.

Let $\mathcal X$ be the feasible states, $\mathcal T\subseteq\mathcal X$
the transversals, and
\begin{equation}\label{eq:graph-masses}
 Z_0=\sum_{X\in\mathcal T}f(X),\qquad
 \Lambda=\sum_{X\in\mathcal X}f(X),\qquad
 \pi(X)=\frac{f(X)}{\Lambda}.
\end{equation}
An original feasible table gives a transversal and a completion by
retaining its small entries and adding $L$ to each large entry.
The small entries respect their capacity because every small cell has
an incident original margin below $U$.  This construction also gives a
specified starting transversal, using the greedy original table.

\subsection{Repairs and graph edges}

For a defect with negative slot $s=(i,j)$ and positive slot
$t=(i',j')$, put $v=(i',j)$.  With $e_w$ denoting a small-slot unit
vector, define its repair $\mathcal R(X)$ by
\begin{equation}\label{eq:repair-map}
 x'=x-e_t+\one_{\{v\in\mathcal S\}}e_v,
 \qquad q'=x',\qquad y'=U\one-q'.
\end{equation}
The term involving $e_v$ is omitted if $v$ is large.
The undirected graph on $\mathcal X$ contains all pairs that differ by
lowering one doubled small coordinate by one and raising another by
one, and also every pair $\{X,\mathcal R(X)\}$ with $X$ a defect.
Repeated edges are counted only once.  Its unnormalized energy is
\begin{equation}\label{eq:graph-energy}
 E(H)=\sum_{\{X,Y\}\text{ an edge}}
 \min\{f(X),f(Y)\}(H(X)-H(Y))^2.
\end{equation}

\begin{lemma}[Repair and preservation of prefixes]\label{lem:repair}
The repair is a feasible transversal.  For each fixed ordered pair
$(s,t)$ it is injective on small states and extends to an injection
\[
 (X,A)\longmapsto
 \begin{cases}
 (\mathcal R(X),A),&v\in\mathcal S,\\
 (\mathcal R(X),A+e_v),&v\in\mathcal B
 \end{cases}
 \quad(A\in\mathcal F_X).
\]
Thus $f(X)\le f(\mathcal R(X))$.
Order small slots by row and then by column.  If $t$ follows $s$, the
repair leaves every slot preceding $s$ unchanged.  When additionally
$x_s=l-1,q_s=l$, its repaired value at $s$ is $l$ if $i'=i$ and $l-1$
otherwise.
\end{lemma}
\begin{proof}
The defect relation is $q=x-e_t+e_s$, and $x_t\ge1$.
If $v$ is small, the changes to $x$ lie in row $i'$, with zero total;
the changes to $q$ are $q'-q=e_v-e_s$, in column $j$, also with zero
total.  Hence all residual margins are unchanged.
The new common entries are nonnegative.  Their row totals are the old
$x$ row totals and their column totals are the old $q$ column totals.
At each small slot, one of these totals is bounded by an original
margin strictly below $U$, proving the capacity constraint.
If $j'=j$, then $v=t$ and the two changes to $x$ cancel; the reset of
$q$ is confined to that column.  If $i'=i$, then $v=s$: the transfer
increments the negative slot, decrements the positive slot, and leaves
$q$ unchanged.  These include both possible coincidences.

If $v$ is large, it differs from $s,t$, and
$x'=x-e_t$, $q'=q-e_s$.  These vectors remain nonnegative and within
capacity, since $x_t\ge1$ and $q_s=x_s+1\ge1$.
Only the residuals of row $i'$ and column $j$ increase, each by one.
Both lines meet the large block.  Adding one to completion entry $v$
therefore gives a valid target completion.  Its entries remain below
$C$ by the total bound already proved.

For fixed $s,t$, the repaired vector recovers
$x=x'+e_t-\one_{\{v\in\mathcal S\}}e_v$, followed by
$q=x-e_t+e_s$ and $y=U\one-q$.  In the large-receiver case subtract
$e_v$ from the target completion; otherwise retain it.  These inverses
on the image prove both injections.

Finally, $t>s$ implies either $i'=i,j'>j$, in which case $v=s$, or
$i'>i$, in which case every small receiver $v$ follows $s$.
Only $s,t,v$ can change, proving prefix preservation and the claimed
special value.
\end{proof}

For later use, fix a balanced prefix $\sigma$ before $s$, and let
$z_-,z_+$ be the total transversal weights extending it with special
values $x_s=l-1,l$, respectively.  Let $c_{st,l}^{\sigma}$ be the total
weight of defects extending that prefix with negative $s$ satisfying
$x_s=l-1,q_s=l$ and positive $t>s$.  Lemma~\ref{lem:repair} maps this
entire class injectively, with completions, into one fixed child.  Hence
\begin{equation}\label{eq:conditional-repair}
 c_{st,l}^{\sigma}\le
 \begin{cases}z_+,&i'=i,\\z_-,&i'\ne i,\end{cases}
 \qquad\text{and in particular}\qquad
 c_{st,l}^{\sigma}\le\max\{z_-,z_+\}.
\end{equation}
Figure~\ref{fig:repair} illustrates the large-receiver case.
\begin{figure}[tb]
\centering
\begin{tikzpicture}[x=1cm,y=1cm,font=\small,
 cell/.style={minimum width=2.4cm,minimum height=1.12cm,inner sep=1pt},
 every path/.style={line width=.55pt}]
 \begin{scope}
  \fill[blue!9] (2.4,-2.24) rectangle (4.8,-1.12);
  \draw (0,0) rectangle (4.8,-2.24);
  \draw (2.4,0)--(2.4,-2.24);
  \draw (0,-1.12)--(4.8,-1.12);
  \node at (2.4,.88) {\textbf{Defect}};
  \node at (1.2,.3) {$c_{j'}<U$};
  \node at (3.6,.3) {$c_j\ge U$};
  \node[anchor=east] at (-.15,-.56) {$r_i<U$};
  \node[anchor=east] at (-.15,-1.68) {$r_{i'}\ge U$};
  \node[cell] at (1.2,-.56) {unchanged};
  \node[cell,align=center] at (3.6,-.56) {$s:\ (a,a+1)$\\[-1pt]\scriptsize negative};
  \node[cell,align=center] at (1.2,-1.68) {$t:\ (b+1,b)$\\[-1pt]\scriptsize positive};
  \node[cell,align=center] at (3.6,-1.68) {$z$\\[-1pt]\scriptsize large entry};
 \end{scope}
 \draw[-{Stealth[length=2mm]},thick] (5.0,-1.12)--(5.9,-1.12);
 \begin{scope}[xshift=6.2cm]
  \fill[blue!9] (2.4,-2.24) rectangle (4.8,-1.12);
  \draw (0,0) rectangle (4.8,-2.24);
  \draw (2.4,0)--(2.4,-2.24);
  \draw (0,-1.12)--(4.8,-1.12);
  \node at (2.4,.88) {\textbf{Repaired transversal}};
  \node at (1.2,.3) {$j'$};
  \node at (3.6,.3) {$j$};
  \node[cell] at (1.2,-.56) {unchanged};
  \node[cell] at (3.6,-.56) {$(a,a)$};
  \node[cell] at (1.2,-1.68) {$(b,b)$};
  \node[cell] at (3.6,-1.68) {$z+1$};
 \end{scope}
\end{tikzpicture}
\caption{The repair when the receiving cell $(i',j)$ is large.
Small cells display $(x,q)$, and the shaded cell displays a completion
entry. Decreasing $x_t$ and matching the viewpoints increases the
residual row margin at $i'$ and column margin at $j$ by one; adding
one to $z$ gives an injective map of completions.
The other cells remain unchanged.}
\label{fig:repair}
\end{figure}


There are fewer than $4d^2$ ordered unit exchanges.  A defect has one
designated repair; a transversal has at most $d^2$ inverse repairs,
one for each ordered defect type.  Thus the degree is at most
$5d^2+1$.  Neighbors can be listed with polynomial bit work using
\eqref{eq:repair-map} and its inverse, testing the boxes, patterns and
residuals, and removing duplicates.  No completion count is needed.
On a repair edge the weight in \eqref{eq:graph-energy} equals the
source defect's weight.

\subsection{Stationary mass and unpadding}

The repair bounds defect mass by transversal mass. To recover an original
table from a padded completion, we also need to control the probability
that a large entry is smaller than the padding. The following estimate
applies to unrestricted tables with any prescribed margins.

\begin{lemma}[Small entries]\label{lem:small-entry}
Consider the uniform law on all nonnegative integer tables with prescribed
margins, with at most $m$ rows and $n$ columns and no entry restrictions.
If both margins incident to a cell are at least $a\ge4d$, then, for every
integer $t$ with $0<t\le a/2$,
\[
 \Pp\{\text{the cell's entry is less than }t\}\le\frac{4td^3}{a}.
\]
\end{lemma}
\begin{proof}
Let the marked cell be $(i,j)$ and its entry be $h<t$.
The other entries in its row have sum at least $a/2$, as do the other
entries in its column.  Unless this event is empty, choose deterministically
$j'\ne j$ and $i'\ne i$ whose donor entries are both at least $a/(2d)$.
Such choices exist because each line has fewer than $d$ entries.  If there
is only one row or column, the event is empty instead.
For every $v=1,\ldots,\lfloor a/(2d)\rfloor$, subtract $v$ from the
entries $(i,j')$ and $(i',j)$ and add $v$ at $(i,j)$ and $(i',j')$.
The resulting table is nonnegative with the same margins.  Given this
table and the labels $(i',j',h)$, its marked entry determines $v$, and
reversing the four changes recovers the source.  There are at most $td^2$
possible labels for each output.  Since
$\lfloor a/(2d)\rfloor\ge a/(4d)$, double counting these moves gives
the asserted bound.  No upper entry bound is imposed in this argument.
\end{proof}

\begin{proposition}[Stationary mass and unpadding]\label{prop:stationary-success}
The masses satisfy
\begin{equation}\label{eq:stationary-mass-bounds}
 0<Z_0\le\Lambda\le(1+d^2)Z_0,
 \qquad \pi_{\min}\ge(C+1)^{-2d}.
\end{equation}
Draw $X$ with law $\pi$ and then a uniform $A\in\mathcal F_X$.
Declare success if $X$ is a transversal and every large entry of $A$
is at least $L$.  Subtracting $L$ from these entries gives a uniform
original table conditional on success, and
\begin{equation}\label{eq:unpadding-success}
 \Pp\{\text{success}\}\ge\frac{1}{2(1+d^2)}.
\end{equation}
\end{proposition}
\begin{proof}
For each of at most $|\mathcal S|(|\mathcal S|-1)\le d^2$ defect types,
the joint injection in Lemma~\ref{lem:repair} bounds its total weight
by $Z_0$.  This proves the upper bound on $\Lambda$; the other mass
inequalities follow from inclusion and the greedy starting table.
There are $\Lambda$ joint objects $(X,A)$, each specified by
$2|\mathcal S|+|\mathcal B|\le2mn<2d$ integer coordinates in $[0,C]$.
Thus $\Lambda\le(C+1)^{2d}$, and the positive integer weights $f(X)$
give the minimum-mass bound.

Consider all unrestricted tables with the enlarged margins
$\bar r,\bar c$, and call their number $F$.  Every transversal and
completion gives one such table, so $Z_0\le F$.
Each large cell has both enlarged margins at least $U$.
Lemma~\ref{lem:small-entry} and a union bound show that the proportion
of these enlarged tables with some large entry below $L$ is at most
\[
 |\mathcal B|\frac{4Ld^3}{U}\le\frac{4Ld^4}{U}
 =\frac4{d^4}<\frac12.
\]
Here $L\le U/2$ and $U\ge4d$; an empty block has no bad cells.
Every enlarged table passing this test can be unpadded to an original
table.  Conversely, padding every original table gives exactly one
passing table and a feasible transversal with its completion.
Consequently the successful pairs are in bijection with
$\Omega(r,c)$ and number at least $F/2\ge Z_0/2$.
Under the stated joint law every pair has probability
$(f(X)/\Lambda)(1/f(X))=1/\Lambda$, proving conditional uniformity
and \eqref{eq:unpadding-success}.
\end{proof}

If there are no small slots, $\mathcal X$ has one state and all graph
energies vanish; the completion and unpadding step still applies.
If there is no large block, the completion step is deterministic.
These conventions will also be used for the Markov chains below.
